\chapter{Testare}

\section{Introducere}

Am testat aplicația manual, pe măsură ce am scris-o, mergând pe scenarii.
Problema e că multe fluxuri au nevoie de mai mulți utilizatori: nu poți testa
căutarea, cererile de colaborare sau o conversație cu un singur cont. Ca să nu
înregistrez de fiecare dată zeci de conturi și să trec prin OTP pentru fiecare,
am scris un serviciu care îmi generează datele de test la pornire.

\section{Generarea datelor de test}

Serviciul rulează doar când mediul de execuție este \texttt{Development}, ca să
nu ajungă niciodată pe un server real. La pornire creează conturi de profesori și
elevi cu profilul completat: vârstă, gen, nivel, materii. Nu scriu conturile
direct în baza de date, ci le trec prin \texttt{UserManager}, ca parolele să fie
procesate de \texttt{PasswordHasher} și rolurile Identity să fie atribuite corect
-- altfel conturile de test s-ar comporta diferit de cele reale, iar testul n-ar
mai însemna nimic. Adaug și câteva cereri de colaborare acceptate, cu conversații
între ele.

Înainte să creeze un cont, serviciul verifică după email dacă există deja. Așa pot
porni aplicația de câte ori vreau fără să se dubleze nimic.

\section{Scenarii testate manual}

\begin{itemize}
    \item \textbf{Înregistrare și verificare}: date valide, cod OTP corect și
    incorect, expirarea codului, formular trimis fără reCAPTCHA;
    \item \textbf{Completarea profilului}: cei trei pași pentru ambele roluri,
    revenirea la un pas anterior cu păstrarea datelor introduse;
    \item \textbf{Căutare și filtrare}: filtrare după nume, rol, nivel și materie,
    izolat și combinat, precum și excluderea conturilor de administrator din
    rezultate;
    \item \textbf{Cereri de colaborare}: trimitere, acceptare, respingere, anulare,
    precum și încercarea unei a doua cereri către același profesor;
    \item \textbf{Mesagerie}: livrare instantanee a unui mesaj text, trimiterea
    unui fișier de către un cont cu rol de profesor și verificarea faptului că
    un elev nu are acces la această funcționalitate, inclusiv prin apelarea
    directă a endpoint-ului API, ocolind interfața;
    \item \textbf{Panou de administrare}: promovarea și retrogradarea unui
    administrator, imposibilitatea unui administrator de a se bloca sau șterge
    pe sine, blocarea temporară și permanentă, precum și verificarea faptului că
    un cont blocat nu se mai poate autentifica;
    \item \textbf{Localizare}: comutarea între română și engleză pe fiecare
    pagină, cu păstrarea limbii alese între navigări succesive.
\end{itemize}

\section{Depanarea unor probleme identificate prin testare}

Testând iar și iar aceleași fluxuri am dat peste trei probleme pe care nu le
văzusem scriind codul:

\begin{itemize}
    \item la primul pas al fluxului de completare a profilului, rolul ales
    (\texttt{Wizard.Role}) nu era transmis prin formular la pasul al doilea,
    fiind folosit doar ca să decid ce câmpuri arăt pe ecran. Așa că la salvare
    aplicația scria mereu în nivelul de predare, oricare ar fi fost rolul ales.
    Am rezolvat punând un câmp ascuns care duce rolul mai departe, de la un pas
    la altul;
    \item controllerul responsabil de afișarea listei de colaborări nu era
    protejat cu \texttt{[Authorize]}, cum erau toate celelalte. Dacă intrai pe
    pagină fără să fii autentificat, primeai o eroare în loc să fii trimis la
    pagina de autentificare. Am adăugat atributul care lipsea;
    \item am încercat să intru pe pagini scriind adresa direct în browser, nu
    apăsând butoanele. Așa am văzut că puteam deschide o conversație cu orice
    cont, doar punându-i identificatorul în URL, chiar dacă nu aveam nicio
    colaborare cu el. Din interfață nu ajungi acolo decât din lista de
    colaborări, dar asta nu e o protecție. Am mutat verificarea pe server, și
    pentru afișarea istoricului, și pentru cele două endpoint-uri prin care se
    scrie un mesaj.
\end{itemize}

Toate trei au ieșit la iveală abia când am încercat aplicația din poziția altui
rol decât cel la care mă gândisem când scriam codul, sau când am ocolit
interfața.
